\section{The original theta complex and its arithmetic packet}
\label{sec:arithmetic-input}
The September 12 source chain has advanced from a support dictionary
to an actual characteristic-zero cochain construction. We retain
its complete coefficient square
\[
\begin{array}{ccc}
G(\mathbb Z)&\xrightarrow{G(j)}&G(\mathbb C)\\
p_{\mathbb Z}\downarrow&&\downarrow p_{\mathbb C}\\
\mathbb Z&\xrightarrow{j}&\mathbb C .
\end{array}
\]
In particular \(p_{\mathbb Z}\) has infinite target. The absolute
pointed base is \(\{\tau,1\}\), with \(\tau\) absorbing for
multiplication; it is not obtained by imposing \(e=1\).
The map \(p_R\) sends both \(\tau\) and \(e_R\) to \(0_R\), while
the support map distinguishes them. All of the linear maps below
lift on each given support fibre; their ordinary zero values lift
to that fibre's supported zero.

For clarity, the fixed section called \(R_0\) in the web notes is denoted
\(R_{\mathrm{ref}}\) here. The corrected sequence keeps its name \(R_m\),
including its genuinely different level-zero member. This is an explicit
notation correspondence between the same maps.

\subsection{Source spaces and fixed arithmetic factors}
Let
\[
 V=\{\phi\in\mathcal S(\mathbb R):\phi(-x)=\phi(x),\
          \phi(0)=0,\ \int_{\mathbb R}\phi(x)\,dx=0\},
\]
\[
 \mathscr B=\{F\in C^\infty(\mathbb R_+):
       p_{b,k}(F)=\sup_{x>0}x^b|D^kF(x)|<\infty
          \text{ for all }b\in\mathbb Z,\ k\geq0\},
 \quad D=-x\partial_x.
\]
Use the original Fourier, theta and Mellin transforms
\[
 \widehat\phi(\xi)=\int_{\mathbb R}\phi(x)e^{-2\pi ix\xi}\,dx,\quad
 \Theta\phi(x)=2\sum_{n\geq1}\phi(nx),\quad
 \mathcal MF(s)=\int_0^\infty F(x)x^s\,\frac{dx}{x}.
 \tag{A1}
\]
The seminorms \(p_{b,k}\) define the strong Schwartz topology.
Logarithmic coordinates give completeness by uniform convergence of
every derivative on compact intervals and the weighted bounds.
The conditions defining \(V\) are closed in the even Schwartz space.

\begin{proposition}[The fixed multiplier and injective theta map]
The map \(\Theta:V\to\mathscr B\) is continuous and injective,
\(D\Theta=\Theta D\), and
\[
 \Theta\widehat\phi(x)=x^{-1}\Theta\phi(1/x).
\]
On \(0<\operatorname{Re}s<1\), define
\[
 H_\phi(s)=\frac{2\pi^{s/2}}
             {s(s-1)\Gamma(s/2)}\int_0^\infty\phi(x)x^s\,\frac{dx}{x}.
\]
Then
\[
 \mathcal M\Theta\phi=gH_\phi,\qquad
 g(s)=2\xi(s)=s(s-1)\pi^{-s/2}\Gamma(s/2)\zeta(s).
 \tag{A2}
\]
For the exact original source
\(\phi_*=(4\pi^2x^4-6\pi x^2)e^{-\pi x^2}\),
\[
 H_{\phi_*}=1,\qquad
 H_{P(D)\phi_*}=P,\qquad
 \mathcal M\Theta(P(D)\phi_*)=gP .
 \tag{A3}
\]
\end{proposition}
\begin{proof}
Schwartz estimates give uniform convergence of each differentiated
theta series on \(x\geq\delta>0\) and all bounds at infinity.
Poisson summation and the two source moments give the displayed
Fourier identity; applying the same bounds to \(\widehat\phi\)
gives all bounds at zero, with continuous dependence on Schwartz
seminorms. Differentiating the series proves intertwining.
For \(x>0\), the absolutely convergent divisor rearrangement gives
\[
 \frac12\sum_{n\geq1}\mu(n)\Theta\phi(nx)
   =\sum_{k\geq1}\left(\sum_{n\mid k}\mu(n)\right)\phi(kx)=\phi(x).
\]
Absolute convergence follows from
\(\sum_k d(k)|\phi(kx)|<\infty\).
This proves injectivity.

In a right half-plane absolute integration gives
\(\mathcal M\Theta\phi=2\zeta M_+\phi\).
Both sides continue to the open strip, with the source's zero
moments removing the relevant apparent singularities, which
proves (A2). More explicitly \(M_+\phi\) is holomorphic for
\(\operatorname{Re}s>-2\), since evenness and \(\phi(0)=0\)
give \(\phi(x)=O(x^2)\) at zero; the theta side is entire.
The multiplier defining \(H_\phi\) is holomorphic and nonzero
on the open strip. The Gaussian integral gives
\[
 M_+\phi_*(s)
 =\pi^{-s/2}\{2\Gamma(s/2+2)-3\Gamma(s/2+1)\}
 =\tfrac12s(s-1)\pi^{-s/2}\Gamma(s/2).
\]
The source moments follow from its value at zero and this formula
at \(s=1\). Integration by parts proves
\(H_{D\phi}=sH_\phi\); all boundary terms vanish.
These calculations prove (A3), with both factors of two retained.
\end{proof}

The actual one-leg complex, in degrees zero and one, is
\(C_+=[V\xrightarrow{\Theta}\mathscr B]\).
Let \(Q=\mathscr B/\Theta V\) and \(q:\mathscr B\to Q\).
The two-leg complex is
\[
 C_{+-}=[V\oplus V\xrightarrow{d}\mathscr B],\quad
       d(v,w)=\Theta(v-\widehat w).
 \tag{A4}
\]
Its source operator is \((D,1-D)\).
Its degree-zero cycles are exactly \((\widehat w,w)\);
its degree-one homology is \(Q\). This follows from injectivity
of theta and \(\operatorname{im}d=\Theta V\).
Thus there is a concrete inclusion of \(C_+\) into \(C_{+-}\),
with the extra degree-zero cycle retained through the map
\(w\mapsto(\widehat w,w)\).

\subsection{The exact operator range, including its inverse}
\begin{proposition}[Source and target Euler inverses]
For \(\rho\in\mathbb C\), on the indicated kernel define
\[
 S_\rho F(x)=x^{-\rho}\int_x^\infty F(y)y^{\rho-1}\,dy .
\]
Then
\[
 (D-\rho)\mathscr B=\ker(\mathcal M(\,\cdot\,)(\rho)),\qquad
 (D-\rho)S_\rho=1,\qquad
 \mathcal M S_\rho F(s)=\frac{\mathcal MF(s)}{s-\rho}.
 \tag{A5}
\]
The operator \(D-\rho\) is injective.
For \(0<\operatorname{Re}\rho<1\), the same formula, extended evenly,
gives
\[
 (D-\rho)V=\ker(H_{(\,\cdot\,)}(\rho)).
 \tag{A6}
\]
These inverses are continuous for the displayed source topologies.
\end{proposition}
\begin{proof}
Differentiation proves the right-inverse identity.
For \(F\) with zero Mellin value the upper tail equals the negative
lower tail. In \(x=e^y\) coordinates, with \(\sigma=\operatorname{Re}\rho\),
the upper bound \(C_Ne^{-Ny}\) becomes
\(C_Ne^{-Ny}/(N-\sigma)\) for \(N>\sigma\), and the lower bound
\(C'_Ne^{Ny}\) becomes \(C'_Ne^{Ny}/(N+\sigma)\) for \(N+\sigma>0\).
Differentiating uses
\(\partial_y S_\rho F(e^y)=-\rho S_\rho F(e^y)-F(e^y)\).
These estimates prove every required seminorm bound and continuity.
The homogeneous solution \(cx^{-\rho}\) fails the target decay at
one end unless \(c=0\), proving injectivity and uniqueness.
Integration by parts proves the Mellin identity, with its removable
value at \(\rho\) filled in.

For the source, \(H_\phi(\rho)=0\) implies \(M_+\phi(\rho)=0\),
since its coefficient is nonzero in the open strip.
Near zero the same inverse is
\[
 S_\rho\phi(x)=-\int_0^1u^{\rho-1}\phi(xu)\,du .
\]
Dominated differentiation gives smooth even extension and value zero.
The upper tail gives Schwartz decay. Integration of
\((D-\rho)S_\rho\phi=\phi\) on \(\mathbb R\) gives
\((1-\rho)\int S_\rho\phi=\int\phi=0\).
Thus the second moment is preserved. The same estimates, with
derivatives under the integral at zero, prove continuity on \(V\).
\end{proof}

For any polynomial
\(h(t)=\prod_i(t-\rho_i)^{m_i}\) with distinct centres in the open
strip, iteration of (A5)--(A6) gives
\[
 h(D)\mathscr B=\ker(j_h\mathcal M),\qquad
 h(D)V=\ker(j_hH),\qquad E_h=\mathbb C[t]/(h).
 \tag{A7}
\]
Here \(j_h\) retains the full Taylor jets through the
Chinese-remainder isomorphism. Order of inversion does not affect
the answer, by uniqueness for the injective \(h(D)\).
The source jet map is onto by (A3).
The target jet map is onto: the finitely many functions
\(y^ke^{\rho_i y}/k!\) are independent on any real interval;
their Gram against a smooth nonnegative bump positive on an
interval is invertible. Its inverse produces compactly supported
smooth log functions with any prescribed vector of moments.
Consequently the actual quotient square is
\[
\begin{array}{ccc}
V/h(D)V&\xrightarrow{\Theta}&\mathscr B/h(D)\mathscr B\\
j_hH\downarrow\wr&&\wr\downarrow j_h\mathcal M\\
E_h&\xrightarrow{\times j_hg}&E_h .
\end{array} \tag{A8}
\]

\begin{theorem}[Finite arithmetic kernel, quotient and packet section]
The actual theta sequence gives the exact sequence
\[
 0\longrightarrow Q[h(D)]\longrightarrow E_h
       \xrightarrow{\times j_hg}E_h
       \longrightarrow Q/h(D)Q\longrightarrow0 .
 \tag{A9}
\]
For a nonempty finite packet \(Z\) of actual nontrivial zeros
with their full orders,
put
\[
 h=\prod_{\rho\in Z}(t-\rho)^{m_\rho},\quad d=\deg h,\quad
 E=E_h,\quad v=g/h,\quad \varepsilon=j_h(h/g).
\]
The element \(\varepsilon\) is a unit. Let \(R_Z(u)\) be the unique
polynomial of degree \(<d\) representing \(\varepsilon u\), and set
\[
 R_{\mathrm{ref}}u=S_h^{\mathscr B}\Theta(R_Z(u)(D)\phi_*),\quad
 J_ZF=j_h\mathcal MF,\quad \sigma=qR_{\mathrm{ref}} .
\]
Then
\[
 \mathcal M R_{\mathrm{ref}}u=(g/h)R_Z(u),\quad J_ZR_{\mathrm{ref}}=1,\quad J_Z\Theta=0,
 \quad Q=Q[h(D)]\oplus h(D)Q .
 \tag{A10}
\]
The projector is \(\sigma\overline J_Z\); its image is \(Q[h(D)]\),
and its kernel is \(h(D)Q\).
\end{theorem}
\begin{proof}
If \(h(D)[F]=0\), write the unique
\(h(D)F=\Theta\phi\) and send \([F]\) to \(j_hH_\phi\).
Changing \(F\) by a theta relation changes \(\phi\) by \(h(D)V\),
so this is well-defined. If its value is zero, (A7) gives
\(\phi=h(D)\psi\), and injectivity yields \(F=\Theta\psi\).
Its image lies in \(\ker(j_hg)\) by (A8).
Conversely a polynomial representative \(P\) of a vector in that
kernel gives the class \([S_h\Theta(P(D)\phi_*)]\).
It is defined by (A7), and changing \(P\) by \(hT\) changes the lift
by \(\Theta(T(D)\phi_*)\). These are inverse maps.
The cokernel description in (A9) follows by quotienting the
right side of (A8), proving exactness with its original arrows.

For the packet, \(j_hg=0\), and \(v\) is nonvanishing at every
chosen zero because the complete order was retained. The Mellin
formula proves \(J_ZR_{\mathrm{ref}}=j_hv\,\varepsilon=1\).
The kernel-to-quotient comparison of (A9) is multiplication by
\(j_hv\), an isomorphism. Its inverse specifies \(\sigma\).
Thus every class has a unique sum of a class in \(Q[h(D)]\)
and one in \(h(D)Q\), proving (A10) and the projector assertions.
\end{proof}

\subsection{The original extension supplies the boundary coefficient}
With \(A=M_t\) on \(E\), define
\(\ell(u)=[t^{d-1}]R_Z(u)\).
Polynomial division, without changing \(R_Z\) or its unit, gives
\[
 DR_{\mathrm{ref}}-R_{\mathrm{ref}}A=f_0\ell,\qquad f_0=\Theta\phi_* .
 \tag{A11}
\]
Indeed \(tR_Z(u)-R_Z(tu)\) is divisible by the monic \(h\),
and its degree is at most \(d\); its quotient is precisely \(\ell(u)\).
Apply (A3) and Mellin injectivity, or multiply by the injective
\(h(D)\), to obtain (A11).
For one simple zero this coefficient is \(1/g'(\rho)=1/(2\xi'(\rho))\).

The invariant subspace
\(\mathscr B_Z=\Theta V+R_{\mathrm{ref}}E\) fits into the actual extension
\[
 0\to V\xrightarrow{\Theta}\mathscr B_Z\xrightarrow{J_Z}E\to0.
 \tag{A12}
\]
It splits complex-linearly through \(R_{\mathrm{ref}}\); in these coordinates its
operator is
\[
 \begin{pmatrix}D_V&\phi_*\ell\\0&A\end{pmatrix}.
\]
The extension class, using the resolution
\(\mathbb C[t]\xrightarrow{h}\mathbb C[t]\), is the class of
\(R_Z(1)(D)\phi_*\) in \(V/h(D)V\).
Under (A7) it is exactly \(\varepsilon\neq0\).
This follows by applying \(h(D)\) to the chosen lift \(R_{\mathrm{ref}}(1)\).
It has no equivariant section into \(\mathscr B\), since every
equivariant image of \(E\) would be killed by \(h(D)\), which is
injective there. Its exact replacement is the equivariant roof
\[
 E[-1]\ \xleftarrow{\simeq}\
   [V\to\mathscr B_Z]\ \longrightarrow\ [V\to\mathscr B].
 \tag{A13}
\]
The left arrow is \(J_Z\) in degree one and zero in degree zero;
exactness of (A12) proves it a quasi-isomorphism.
Thus no relation or homotopy is suppressed by the ordinary
direct-sum notation.

\subsection{Actual representative corrections and rank-two control}
Throughout this subsection \(\mathscr H=L^2(\mathbb R_{>0},dx)\),
with inner product conjugate-linear in the first slot. Coefficient
spaces carry the Euclidean inner products in the retained bases;
every adjoint below uses these specified metrics.
Put \(\phi_j=D^j\phi_*\), \(f_j=D^jf_0\), and
\(\Phi_m(c)=\sum_{j=0}^mc_j\phi_j\).
For \(F_m=\Theta\Phi_m:\mathbb C^{m+1}\to\mathscr H\), define
\[
 \begin{gathered}
 H_m=F_m^*F_m,\quad C_m=F_m^*R_{\mathrm{ref}},\quad
 B_m=-H_m^{-1}C_m,\\
 \Pi_m=F_mH_m^{-1}F_m^*,\quad
 R_m=R_{\mathrm{ref}}+F_mB_m=(1-\Pi_m)R_{\mathrm{ref}}.
 \end{gathered}
 \tag{A14}
\]
The columns \(f_j\) are independent: Mellin would turn a dependence
into \(gP=0\). Thus \(H_m>0\).
The correction is an original theta boundary, giving
\(qR_m=\sigma\), \(J_ZR_m=1\), and \(F_m^*R_m=0\).
For every original coefficient matrix \(B:E\to\mathbb C^{m+1}\),
\[
 \begin{gathered}
 (R_{\mathrm{ref}}+F_mB)^*(R_{\mathrm{ref}}+F_mB)
       =G_m+(B-B_m)^*H_m(B-B_m),\\
 G_m=R_m^*R_m=R_{\mathrm{ref}}^*R_{\mathrm{ref}}-C_m^*H_m^{-1}C_m>0 .
 \end{gathered}
 \tag{A15}
\]
Expansion proves the identity; injectivity from \(J_ZR_m=1\)
proves its strict positivity.

Let \(I_me_j=e_j\) and \(S_me_j=e_{j+1}\), both with codomain
\(\mathbb C^{m+2}\). In the retained source coordinates,
\[
 K_m=e_0\ell+S_mB_m-I_mB_mA,\qquad
 DR_m-R_mA=F_{m+1}K_m .
 \tag{A16}
\]
Integration by parts in \(dx\), with the source boundary term
\([x\overline FH]_0^\infty=0\), gives
\[
 W_m=A^*G_m+G_mA-G_m
       =-(R_m^*F_{m+1}K_m+K_m^*F_{m+1}^*R_m).
\]
All but the last source column pair to zero.
With \(z_m=R_m^*f_{m+1}\) and \(b_m=e_m^*B_m\) this is exactly
\[
                   W_m=-z_mb_m-b_m^*z_m^* .
 \tag{A17}
\]
In particular \(\operatorname{rank}W_m\leq2\).
Factoring \(G_m^{-1}W_m\) through the two columns gives its nonzero
spectrum as that of
\[
 -\begin{pmatrix}c_m&v_m\\u_m&\overline c_m\end{pmatrix},\quad
 u_m=z_m^*G_m^{-1}z_m,\quad v_m=b_mG_m^{-1}b_m^*,\quad
 c_m=b_mG_m^{-1}z_m.
\]
Cauchy--Schwarz gives \(u_mv_m\geq|c_m|^2\).
Hence its exact two-sided allowance is
\[
 \epsilon_m=|\operatorname{Re}c_m|+
            \sqrt{u_mv_m-(\operatorname{Im}c_m)^2},\qquad
 -\epsilon_mG_m\leq W_m\leq\epsilon_mG_m .
 \tag{A18}
\]
This includes ranks zero and one.
For a reflection-stable packet,
\(\operatorname{Tr}(G_m^{-1}W_m)
=2\operatorname{Re}\operatorname{Tr}A-d=0\).
Every eigenvector retains the equality
\[
        \frac{u^*W_mu}{u^*G_mu}=2\operatorname{Re}\rho-1 .
 \tag{A19}
\]
The next sections compute the actual arithmetic moment problem
behind \(G_m\), sharpen its constraints, and determine the topology
of the original jet observation. They use (A11)--(A19) rather than
an assigned metric in place of the theta representative.

